%!TEX output_directory = .cache

\documentclass[12pt]{letter}
\usepackage[utf8]{inputenc}
\usepackage{graphicx}
\usepackage{eurosym}
\usepackage{charter}
\usepackage{hyperref}
\usepackage{xcolor}

\hypersetup{
    colorlinks,
    linkcolor=[HTML]{404040},
    urlcolor={red!50!black}
}

\signature{\includegraphics[scale=0.6]{images/alessandro}
			Alessandro Tommasin \\
			\textit{Responsabile di progetto}
			\href{mailto:redroundrobin.site@gmail.com}{redroundrobin.site@gmail.com}}
\address{ Via Trieste 63 \\ Padova \\ PD 35121, Italia}

\date{13 aprile 2020}

\begin{document}

\begin{letter}{ }

\includegraphics[scale=0.17]{images/logo.png}

\opening{Gentile prof. Vardanega,\\ Gentile prof. Cardin, }

con la presente il gruppo 6, \textit{Red Round Robin}, vorrebbe richiedere la partecipazione alla \textit{Revisione di Qualifica} in qualità di fornitori per il progetto \textbf{ThiReMa (Capitolato 6)}, proposto dall'azienda \textit{SanMarco Informatica}.
\newline
Inoltre, comunichiamo che le regole di versionamento dei documenti hanno subito alcune variazioni, come riportato nelle \textit{Norme di Progetto v2.3.0+b0.15}; questo ha comportato una cambio di versione di tutti i documenti.

Le forniamo il link di una repository pubblica da cui è possibile scaricare e visualizzare online i documenti richiesti dal progetto.

In particolare:

\begin{itemize}
  \item \textit{Analisi dei Requisiti v2.5.0+b0.15};
	\item \textit{Piano di Qualifica v1.6.0+b0.15};
	\item \textit{Piano di Progetto v2.5.0+b0.15};
	\item \textit{Norme di Progetto v2.3.0+b0.15};
	\item \textit{Glossario v1.4.0+b0.15};
	\item \textit{Manuale del manutentore v1.0.0+b0.15};
	\item \textit{Manuale dell'utente v1.0.0+b0.15};
	\item \textit{verbali interni} dal 2020-03-09 a oggi;
	\item \textit{verbali esterni} dal 2020-03-09 a oggi.
\end{itemize}

Il link per scaricare e visualizzare i documenti è il seguente:

\begin{center}
\href{}{http://rq.redroundrobin.site}
\end{center}

%\newpage

Per qualunque altro chiarimento o problema, rimaniamo a completa disposizione.

\closing{Cordiali saluti,}


\vspace{3em}
\ps

\textbf{P.S.} Se il link fornito non dovesse funzionare, potete fare riferimento al seguente link:
\href{https://www.redroundrobin.site}{www.redroundrobin.site}

\end{letter}

\end{document}
